\chapter{Databases and SQL}\label{ch:pl:databases-and-sql}

A regulator asks for the firm's positions report exactly as it stood at 18:00 last Tuesday. On Wednesday morning operations had corrected
Monday's and Tuesday's errors --- quantities keyed ten times too large, trades that never happened, trades booked a day late --- and they
had corrected them in place. The positions table now gives the right answer for Tuesday evening, and nobody can say what answer the firm gave
on Tuesday evening. This chapter builds the firm's trade database so that both answers stay available: relational tables with keys, writes in
transactions that do not lose updates, indexes chosen by reading \omterm{def:pl:databases-and-sql:index}{query plans}, and a second time axis that records what the database believed
and when.

\section{The relational model for trades and positions}

\begin{definition}[Relational model, primary key, foreign key]\label{def:pl:databases-and-sql:relational}
The \emph{relational model}\index{relational model}, proposed by Codd in 1970, stores data as relations --- tables of rows with named, typed
columns --- manipulated by operations on whole relations (selection, projection, join), independently of how rows are stored. A \emph{primary
key}\index{primary key} is a set of columns whose values identify one row; a \emph{foreign key}\index{foreign key} is a set of columns whose
values must appear as the primary key of a row in another table.
\end{definition}

\begin{definition}[Normal form]\label{def:pl:databases-and-sql:normal}
A schema is in a \emph{normal form}\index{normal form} when each fact is stored once: every non-key column depends on the key, the whole key
and nothing but the key (third normal form), so that an instrument's currency, say, lives in the instruments table and not on every trade.
\end{definition}

The chapter's schema (\texttt{firm.tradedb}) has three tables: instruments, accounts and trades, each trade pointing to its account and its
instrument by \omterm{def:pl:databases-and-sql:relational}{foreign keys} that the database enforces --- a trade for an account that does not exist is refused, not stored. Positions are
not a table: they are a query over trades, so they cannot disagree with them. Schema changes are numbered migrations, applied once each and
recorded in the database, so that every copy of the database can say which schema it has.

\omcode{code/firm/tradedb/firm_tradedb.py}{26}{41}{The schema as migrations: primary keys, foreign keys, and a trades table with two time
intervals per row.}\label{lst:pl:databases-and-sql:schema}

\begin{omfigure}
\begin{tikzpicture}[font=\small,
  tab/.style={draw=omDef,fill=omDef!6,rounded corners=2pt,align=left,text width=38mm,inner sep=4pt},
  arr/.style={-{Stealth[length=2mm]},thick,omThm}]
\node[tab] (i) at (0,1.4) {\textbf{instruments}\\\underline{id}, symbol (unique), currency};
\node[tab] (a) at (0,-0.9) {\textbf{accounts}\\\underline{id}, name (unique)};
\node[tab,text width=64mm] (t) at (7.4,0.25) {\textbf{trades}\\\underline{trade\_id}, account\_id, instrument\_id,\\qty, price,
valid\_from, valid\_to,\\\underline{sys\_from}, sys\_to};
\draw[arr] (t.west |- 0,0.75) -- node[above,font=\scriptsize]{instrument\_id} (i.east |- 0,0.75);
\draw[arr] (t.west |- 0,-0.55) -- node[below,font=\scriptsize]{account\_id} (a.east |- 0,-0.55);
\node[font=\scriptsize,text width=64mm,align=left] at (7.4,-1.75) {primary key underlined; arrows are foreign keys};
\end{tikzpicture}

\omcaption{The chapter's schema. Instruments and accounts are referenced by trades through enforced \omterm{def:pl:databases-and-sql:relational}{foreign keys}; the trades table's \omterm{def:pl:databases-and-sql:relational}{primary
key} is the trade identifier with the start of its \omterm{def:pl:databases-and-sql:system}{system time}, because a corrected trade has several versions.}
\label{fig:pl:databases-and-sql:schema}
\end{omfigure}

\section{Transactions and isolation}

\begin{definition}[ACID transaction]\label{def:pl:databases-and-sql:acid}
An \emph{ACID transaction}\index{ACID transaction} is a group of reads and writes that is atomic (all or nothing), consistent (it takes the
database from one valid state to another), isolated (concurrent transactions do not see each other's partial work) and durable (once committed,
it survives a crash).
\end{definition}

\begin{definition}[Isolation level, write-ahead log]\label{def:pl:databases-and-sql:isolation}
An \emph{isolation level}\index{isolation level} states which anomalies of concurrency a database allows between transactions, from read
uncommitted to serializable, under which the result is that of some serial order. A \emph{write-ahead log}\index{write-ahead log} records each
change before it is applied to the database's main file, so that a crash can be recovered by replaying or discarding it: in SQLite's WAL mode
\textquotedbl{}the original content is preserved in the database file and the changes are appended into a separate WAL file\textquotedbl{}.
\end{definition}

The anomaly that costs a trading firm money is the lost update, which Berenson and his co-authors define as a transaction reading an item,
another updating it, and the first then writing a value computed from its earlier read, so that the second's update is lost. Two gateways
each add 100 shares to the same position: each reads the position, adds its fill, and writes the result. SQLite is serializable --- it
\textquotedbl{}implements serializable transactions by actually serializing the writes\textquotedbl{} --- but only within a transaction; an
application that reads in one statement and writes in another has made two transactions, and the database cannot protect it.

\omcode{code/platforms/25-databases-and-sql/python/pl_tradedb.py}{132}{144}{One writer's two steps in each of three modes: read then write
as two transactions, the same inside one transaction taken at BEGIN, or one statement that reads and writes.}
\label{lst:pl:databases-and-sql:lost}

Over 1\,000 random interleavings of the two writers' steps, the read-then-write application loses an update 636 times, close to the two
thirds of interleavings in which both reads come before a write. Inside a transaction taken at \texttt{BEGIN IMMEDIATE}, the second writer is
refused until the first commits and then retries: no update is lost. A single \texttt{UPDATE pos SET q = q + 100} does the read and the write
in one statement and loses nothing either. The rule: send an increment to the database as an increment, or let the read and the write
share one transaction, with a retry.

\section{Indexes and query plans}

\begin{definition}[B-tree index, query plan]\label{def:pl:databases-and-sql:index}
A \emph{B-tree index}\index{B-tree index} keeps a copy of chosen columns of a table in a balanced sorted tree with pointers to the rows, so
that rows with given values or in a range are found in logarithmic time instead of by reading the table. A \emph{query plan}\index{query
plan} is the sequence of operations the database chooses to answer a query --- which tables it scans, which indexes it searches, how it joins
and groups --- and can be displayed before the query runs.
\end{definition}

Three queries run over the week's 500\,008 trades: a point query (one account's position in one instrument as of Tuesday 18:00), the
report (every account's positions as of Tuesday 18:00, joined to the account and instrument names), and an analytical query (the week's
traded value by instrument). \texttt{EXPLAIN QUERY PLAN} shows each one reading the whole trades table (\texttt{SCAN trades}); the chapter's
advisor lists such scans, and the question for each is whether an index would select few enough rows to be worth it.

\begin{omfigure}
\begin{tikzpicture}
\begin{semilogyaxis}[width=11.5cm,height=5.6cm,ybar,bar width=8pt,ylabel={milliseconds (median of 5, log scale)},symbolic x coords={point,report,analytic},
  xtick=data,xticklabels={point query,positions report,traded value by instrument},enlarge x limits=0.22,
  grid=major,grid style={gray!20},tick label style={font=\small},label style={font=\small},table/col sep=comma,ymin=0.005,ymax=2000,
  log origin=infty,
  legend cell align=left,
  legend style={font=\scriptsize,at={(0.5,1.03)},anchor=south,legend columns=2,draw=none,/tikz/every even column/.append style={column sep=3mm}}]
\addplot[fill=omThm!60,draw=omThm] table[x=query,y=sqlite_none]{figdata/platforms/25-databases-and-sql/queries_wide.csv};
\addplot[fill=omDef!60,draw=omDef] table[x=query,y=sqlite_ai]{figdata/platforms/25-databases-and-sql/queries_wide.csv};
\addplot[fill=omDef!20,draw=omDef] table[x=query,y=sqlite_v]{figdata/platforms/25-databases-and-sql/queries_wide.csv};
\addplot[fill=gray!50,draw=gray] table[x=query,y=duckdb]{figdata/platforms/25-databases-and-sql/queries_wide.csv};
\legend{SQLite without index,SQLite with index on (account; instrument),SQLite with index on valid time,DuckDB (no index)}
\end{semilogyaxis}
\end{tikzpicture}

\omcaption{Measured query times over 500\,008 trades (log scale; one core, in memory, median of five runs, on an otherwise
idle machine). The index on (account, instrument) takes the point query from 13.7~ms to 0.01~ms; the index on valid time barely helps any query,
because Tuesday 18:00 still selects two days of five; DuckDB runs the report 15 times and the analytical query 53 times faster than SQLite
without any index. Data: \texttt{bench\_tradedb.py}, \texttt{measured\_queries.csv}.}\label{fig:pl:databases-and-sql:queries}
\end{omfigure}

The measurements (\cref{fig:pl:databases-and-sql:queries}) give the rule of thumb its numbers. An index pays when it selects a small part of
the table: the point query reads about fifty rows of 500\,008 through the (account, instrument) index, a gain of three orders of magnitude.
An index on valid time selects 40\% of the rows for Tuesday's report and saves almost nothing, and no index helps a query that must read every
row. Each index also costs every insert its own write, so the advisor's list is a list of candidates, not of indexes to create.

\section{Time in the database: system time and time travel}

Book~7, chapter~3, stored data with two times, the valid time a value describes and the knowledge time at which it became available, and made
research reproducible with them. A trade database needs the same two axes for a slightly different reason.

\begin{definition}[System time, time-travel query]\label{def:pl:databases-and-sql:system}
The \emph{system time}\index{system time} of a row version is the interval during which the database held it as current: from the
transaction that inserted it to the transaction that superseded it. A \emph{time-travel query}\index{time-travel query} asks for the rows
valid at one time as the database held them at another, by filtering on both intervals.
\end{definition}

\omterm{def:pl:databases-and-sql:system}{System time} differs from Book~7's knowledge time in one respect: it is when the firm's database recorded a fact, which the database stamps
itself and nobody can backdate, where knowledge time is when the fact became available, which may be earlier and is supplied by the data.
SQL:2011 standardised both axes: transaction time through system-versioned tables, valid time through application-time periods, and queries
of the form \texttt{FOR SYSTEM\_TIME AS OF}.

\omcode{code/firm/tradedb/firm_tradedb.py}{82}{92}{A correction closes the old version's system time and inserts the corrected version:
nothing that was believed is overwritten.}\label{lst:pl:databases-and-sql:correct}

\begin{omfigure}
\begin{tikzpicture}[x=0.85cm,y=0.85cm,font=\small]
\draw[-{Stealth[length=2mm]}] (0,0) -- (12.2,0) node[right,font=\scriptsize]{system time};
\draw[-{Stealth[length=2mm]}] (0,0) -- (0,4.4) node[above,font=\scriptsize,align=center]{valid time};
\foreach \d/\n in {0/Mon,2.4/Tue,4.8/Wed,7.2/Thu,9.6/Fri} {\draw (\d,0) -- (\d,-0.1) node[below,font=\scriptsize]{\n};
  \draw (0,\d*0.4) -- (-0.1,\d*0.4) node[left,font=\scriptsize]{\n};}
\fill[omThm!25,draw=omThm] (1.2,0.48) rectangle (5.8,4.1);
\fill[omDef!25,draw=omDef] (5.8,0.48) rectangle (12,4.1);
\node[align=center,font=\scriptsize] at (2.7,3.0) {version 1:\\700 shares\\known from\\Monday 12:01};
\node[align=center,font=\scriptsize] at (8.9,3.3) {version 2: 70 shares\\known from Wednesday 10:00};
\draw[black,thick,dashed] (4.2,0) -- (4.2,4.2);
\draw[black,thick,dashed] (0,1.68) -- (12,1.68);
\fill (4.2,1.68) circle (2.2pt) node[below right,font=\scriptsize,fill=white,inner sep=1pt]{report as it was: Tue 18:00 / Tue 18:00};
\fill (11.6,1.68) circle (2.2pt) node[above left,font=\scriptsize,fill=white,inner sep=1pt]{as corrected: Tue 18:00 / Fri 20:00};
\end{tikzpicture}

\omcaption{One corrected trade in the two times: executed on Monday at 12:00, keyed as 700 shares, corrected to 70 on Wednesday at 10:00.
Each version is a rectangle, true from Monday in valid time (up) and current for an interval of \omterm{def:pl:databases-and-sql:system}{system time} (across). Tuesday's report as it
was reads the point (valid Tuesday 18:00, system Tuesday 18:00) and finds 700; the corrected report reads (Tuesday 18:00, Friday) and finds
70.}\label{fig:pl:databases-and-sql:bitemporal}
\end{omfigure}

\omcode{code/firm/tradedb/firm_tradedb.py}{102}{118}{Time travel: the trades valid at \texttt{:v} as held at \texttt{:s}, and the
positions computed from them.}\label{lst:pl:databases-and-sql:asof}

The week of the model has 500\,008 trades. On Wednesday at 10:00 operations correct 30 quantities keyed ten times too large, cancel 10 trades
that never happened, and book 8 of Tuesday's trades that had been missed; the table then holds 500\,048 row versions, forty more than trades.
Tuesday's 18:00 report, as it was, reads valid time Tuesday 18:00 at \omterm{def:pl:databases-and-sql:system}{system time} Tuesday 18:00; the corrected report reads the same valid time
at \omterm{def:pl:databases-and-sql:system}{system time} Friday evening. The two differ on 48 of the 10\,000 account and instrument positions --- one per correction
(\cref{fig:pl:databases-and-sql:diff}) --- and the regulator gets the first, with the second and the audit table explaining every difference.
The table corrected in place can only produce the second.

\begin{omfigure}
\begin{tikzpicture}
\begin{axis}[width=11.5cm,height=5.4cm,xlabel={the 48 positions that differ, sorted},ylabel={corrected minus as reported (shares)},
  grid=major,grid style={gray!20},tick label style={font=\small},label style={font=\small},table/col sep=comma,
  xmin=0,xmax=49,ymin=-1000,ymax=1100,only marks,yticklabel style={/pgf/number format/1000 sep={\,}},
  legend style={font=\small,at={(0.02,0.98)},anchor=north west,draw=none,fill=white}]
\addplot[scatter,only marks,mark=*,mark size=1.8pt,point meta=explicit symbolic,
  scatter/classes={qty={omThm},cancel={black},late={omDef}}]
  table[x=rank,y=change,meta=kind]{figdata/platforms/25-databases-and-sql/report_diff.csv};
\node[omThm,font=\scriptsize,anchor=west] at (axis cs:2,920) {$\bullet$ quantity corrected (30)};
\node[font=\scriptsize,anchor=west] at (axis cs:2,760) {$\bullet$ trade cancelled (10)};
\node[omDef,font=\scriptsize,anchor=west] at (axis cs:2,600) {$\bullet$ late trade booked (8)};
\end{axis}
\end{tikzpicture}

\omcaption{Tuesday's 18:00 positions report, corrected minus as it was reported, for the 48 account and instrument positions that differ:
quantity corrections (red) take nine tenths of the trade out of the position, cancellations (black) the whole trade, and late bookings
(blue) add one. The other 9\,952 positions agree. Data: \texttt{fig\_tradedb.py}.}\label{fig:pl:databases-and-sql:diff}
\end{omfigure}

\section{Transactional and analytical workloads}

\begin{definition}[Online transaction processing, online analytical processing]\label{def:pl:databases-and-sql:oltp}
\emph{Online transaction processing}\index{online transaction processing} (OLTP) is a workload of many small transactions that read and write
a few rows each by key --- booking trades, updating positions. \emph{Online analytical processing}\index{online analytical processing} (OLAP)
is a workload of few large queries that read many rows and few columns to aggregate them --- traded value by instrument, P\&L by desk over a
year.
\end{definition}

SQLite is a row store built for the first; DuckDB, the \omterm{def:pl:time-series-databases-and-the-alternatives:embedded}{embedded database} of chapter~\ref{ch:pl:time-series-databases-and-the-alternatives},
is a columnar engine with vectorised execution built for the second. On the same SQL and the same data --- the chapter checks that every
answer is equal --- DuckDB runs the report in 9.8~ms against SQLite's 150 and the analytical query in 2.8~ms against 148, with one thread and no
index. SQLite with the right index still wins the point query by far (0.01~ms against 1.0). The firm keeps its trades in a transactional
store, where the keys, the transactions and the \omterm{def:pl:databases-and-sql:system}{system time} are, and gives research and reporting a columnar copy, refreshed from the
transactional one and never written by anyone else.

\begin{dated}{2026-09}{Isolation in SQLite and temporal SQL}\label{dat:pl:databases-and-sql:temporal}
SQLite's documentation states that its transactions are serializable, implemented by serializing the writes (except shared-cache connections
with \texttt{read\_uncommitted}), and that in write-ahead-log mode the changes are appended to a separate WAL file while the database file keeps
the original content. The SQL:2011 standard provides transaction time through system-versioned tables and valid time through application-time
period tables, as described by Kulkarni and Michels (2012).
\end{dated}

\section{Tutorial: Tuesday at 18:00}\label{tut:pl:databases-and-sql:tut}

\begin{tutorial}
\textbf{Goal.} Store a week of trades bitemporally, correct them, reproduce a past report, and measure what indexes and engines change.
\textbf{End state:} \cref{fig:pl:databases-and-sql:diff,fig:pl:databases-and-sql:queries}.
\begin{tutsteps}
\item \textbf{Schema}: \texttt{firm\_tradedb.connect()} and \texttt{migrate}; try inserting a trade for an unknown account.
\item \textbf{The week}: \texttt{pl\_tradedb.load(conn)}; the corrections of Wednesday 10:00.
\item \textbf{Reports}: \texttt{positions(conn, TUE\_18, TUE\_18)} and \texttt{positions(conn, TUE\_18, NOW)}; the rows that differ.
\item \textbf{Lost updates}: \texttt{lost\_updates(mode)} for the three modes.
\item \textbf{Plans and times}: \texttt{plan} and \texttt{advise} for each query; \texttt{bench\_tradedb.py}.
\end{tutsteps}
\textbf{What to change next.} Add an index on (valid\_from, account\_id) and read the report's plan again; copy the tables to DuckDB with a
system-time filter so that the columnar copy can travel in time too.
\end{tutorial}

\section{Build: the trade database}\label{bld:pl:databases-and-sql:tradedb}

\begin{build}
\bfield{Purpose} One transactional store of the firm's trades that enforces its keys, never loses an update, and can reproduce any report as
it was produced.
\bfield{Interface} \texttt{connect}, \texttt{migrate}, \texttt{book}, \texttt{correct}, \texttt{cancel}, \texttt{as\_of},
\texttt{positions}, \texttt{transact}, \texttt{plan}, \texttt{advise}, \texttt{to\_duckdb}.
\bfield{Rules} \omterm{def:pl:databases-and-sql:relational}{Foreign keys} on; schema only through migrations; corrections close and insert, never update values; every correction
audited; read-modify-write only inside a transaction, or as one statement; indexes chosen from plans and measured.
\bfield{Acceptance tests} \texttt{code/firm/tradedb/tests/}: keys enforced; a correction, a cancellation and a late booking seen as they were
and as corrected; a locked writer retried and a failing transaction rolled back; the advisor before and after an index; the same answer in
DuckDB.
\bfield{Stretch} Valid-time corrections that split a row's interval; system-versioned tables in a database that has them; a columnar copy
refreshed by change capture.
\end{build}

\begin{omsources}
\item E.~F.~Codd, \emph{A Relational Model of Data for Large Shared Data Banks}, Communications of the ACM 13(6), 1970.
\item H.~Berenson, P.~Bernstein, J.~Gray, J.~Melton, E.~O'Neil and P.~O'Neil, \emph{A Critique of ANSI SQL Isolation Levels}, SIGMOD 1995.
\item K.~Kulkarni and J.-E.~Michels, \emph{Temporal Features in SQL:2011}, SIGMOD Record 41(3), 2012.
\item SQLite documentation: \emph{Isolation in SQLite}; \emph{Write-Ahead Logging}.
\item One Quant Book~7, chapter~3 (valid time, knowledge time, bitemporal data).
\end{omsources}

\section{Exercises}

\begin{exercise}[$\star$]\label{exo:pl:databases-and-sql:1}
Why are positions a query over trades rather than a table of their own?
\end{exercise}

\begin{exercise}[$\star$]\label{exo:pl:databases-and-sql:2}
Why is the trades table's \omterm{def:pl:databases-and-sql:relational}{primary key} (trade identifier, start of \omterm{def:pl:databases-and-sql:system}{system time}) and not the trade identifier alone?
\end{exercise}

\begin{exercise}[$\star$]\label{exo:pl:databases-and-sql:3}
How many row versions does the table hold after Wednesday, and why forty more than trades?
\end{exercise}

\begin{exercise}[$\star\star$]\label{exo:pl:databases-and-sql:4}
Why does the read-then-write application lose an update in about two thirds of the interleavings?
\end{exercise}

\begin{exercise}[$\star\star$]\label{exo:pl:databases-and-sql:5}
Why does the index on valid time barely help the positions report?
\end{exercise}

\begin{exercise}[$\star\star$]\label{exo:pl:databases-and-sql:6}
A late trade for Tuesday is booked on Wednesday. Which of Tuesday's two reports shows it, and why?
\end{exercise}

\begin{exercise}[$\star\star\star$]\label{exo:pl:databases-and-sql:7}
\emph{Coding.} Write the query that lists every trade whose Tuesday 18:00 version differs from its current version, with both values.
\end{exercise}

\begin{exercise}[$\star\star\star$]\label{exo:pl:databases-and-sql:8}
\emph{Find the flaw.} \textquotedbl{}We keep an audit table of every change, so we can always rebuild what the positions table said on any
day.\textquotedbl{}
\end{exercise}

\section{Problem: Tuesday at 18:00}

\begin{problem}[{Weekend problem --- Tuesday at 18:00}]\label{pb:pl:databases-and-sql:1}
The chapter's week of trades, its corrections, its writers and its queries.

\medskip\noindent\textbf{Part I --- The schema.}
\begin{enumerate}
\item What do the primary and \omterm{def:pl:databases-and-sql:relational}{foreign keys} of the schema guarantee?
\item What does \omterm{def:pl:databases-and-sql:normal}{normal form} keep off the trades table?
\item Why are schema changes migrations?
\item What does each row of the trades table carry beyond the trade?
\item How is a trade corrected, and how is it cancelled?
\end{enumerate}

\medskip\noindent\textbf{Part II --- Concurrency.}
\begin{enumerate}[resume]
\item What is a lost update, in the terms of the 1995 critique?
\item Why does SQLite's serializable isolation not prevent it for the read-then-write application?
\item How many updates are lost in 1\,000 interleavings under each mode?
\item What does \texttt{BEGIN IMMEDIATE} change?
\item When is one statement enough?
\end{enumerate}

\medskip\noindent\textbf{Part III --- Plans and engines.}
\begin{enumerate}[resume]
\item What does the plan of each query show without an index?
\item What does each index do to each query?
\item How do SQLite and DuckDB compare on each query?
\item Why does an index cost something even when it is not used?
\item Which store should research query?
\end{enumerate}

\medskip\noindent\textbf{Part IV --- The verdict.}
\begin{enumerate}[resume]
\item State the \emph{named result}: the positions report reproduced as it stood at 18:00 on the Tuesday against the corrected report, the
rows that differ, and the query-time gain from each index.
\item How does \omterm{def:pl:databases-and-sql:system}{system time} differ from Book~7's knowledge time?
\item What does the table corrected in place lose?
\item What would you ask of any table a report is produced from?
\item In one sentence: what must a database remember besides the right answer?
\end{enumerate}
\end{problem}

\section{Interview questions}

\begin{interviewq}[$\star$ \iqroles{developer}]\label{iq:pl:databases-and-sql:1}
What is a transaction, and what do the four letters of ACID mean?
\end{interviewq}

\begin{interviewq}[$\star\star$ \iqroles{developer}]\label{iq:pl:databases-and-sql:2}
Two services update the same position at the same time. How can an update be lost, and how do you prevent it?
\end{interviewq}

\begin{interviewq}[$\star\star$ \iqroles{developer}]\label{iq:pl:databases-and-sql:3}
A query is slow. How do you decide whether an index will help?
\end{interviewq}

\begin{interviewq}[$\star\star$ \iqroles{developer}]\label{iq:pl:databases-and-sql:4}
A regulator asks for a report exactly as you produced it three months ago. How do you design for that?
\end{interviewq}

\begin{interviewq}[$\star\star\star$ \iqroles{developer}]\label{iq:pl:databases-and-sql:5}
Design the storage of a trading firm's trades and positions for booking, reporting and research.
\end{interviewq}
